% ============================================
% MEM-26 (ZH)
% 中文翻译 — SPEC-935-R356
% ============================================
\section*{\textbf{MEM-26} --- 最后一页}
\markboth{MEM-26 --- 最后一页}{}

\noindent \textit{死前3天}

\noindent \textbf{1979 —— 死前三天}

\vspace{0.5em}\noindent\rule{\textwidth}{0.5pt}\vspace{0.5em}

\noindent 我知道自己快死了。

\noindent 身体正在停止。心脏在衰竭。呼吸变短。怪物在我体内太重了。

\noindent 死前三天，我写下最后一页。

\noindent \textit{“读到这段话的人：你就是下一个。}

\noindent \textit{我背负怪物六十二年。它在我还是个孩子时进入我体内，在一个乱葬坑里，在一座集中营中间。它从里面吃我。但它也让我活了下来。}

\noindent \textit{它需要宿主。我死的时候，它会寻找下一个。}

\noindent \textit{不是一个身体。是一个读者。}

\noindent \textit{你会读这本日记。你读的时候，它会进入你。不是肉体上。是精神上。这个故事会进入你的脑子。而一旦进去，就再也出不来。}

\noindent \textit{你会开始梦见乱葬坑。会闻到气味。会看见房间角落里的黄眼。}

\noindent \textit{你会以为这是暗示。}

\noindent \textit{不是。}

\noindent \textit{是怪物。}

\noindent \textit{欢迎来到循环，我的继任者。}

\noindent \textit{--- 若阿金·安东尼奥·科雷亚}
\textit{--- 患者1,260}
\textit{--- 最后一句话：饥饿。”}\NE{编者记录：本碎片的原始纸页于1979年掘尸时，在若阿金外套的内袋中被发现，是散落的一页。纸张沾满泥土与潮气；墨迹在多处洇开，尤其是下缘。此处转录的是修复后清晰可读的文本。最后一句话——‘饥饿’——几乎无法辨认，仅靠钢笔在纸上的压痕保存下来，即使在墨迹溶解之后，压痕仍在下一页上留下了印记。}

\vspace{0.5em}
\noindent\textit{\textbf{↪ 路线:} \rota{DOC-06} \rotasep \rota{DOC-08}}
